
\documentclass{article}
\usepackage{colortbl}
\usepackage{makecell}
\usepackage{multirow}
\usepackage{supertabular}

\begin{document}

\newcounter{utterance}

\centering \large Interaction Transcript for game `clean\_up', experiment `1\_medium\_7obj\_cs', episode 0 with glm{-}5.2.
\vspace{24pt}

{ \footnotesize  \setcounter{utterance}{1}
\setlength{\tabcolsep}{0pt}
\begin{supertabular}{c@{$\;$}|p{.15\linewidth}@{}p{.15\linewidth}p{.15\linewidth}p{.15\linewidth}p{.15\linewidth}p{.15\linewidth}}
    \# & $\;$A & \multicolumn{4}{c}{Game Master} & $\;\:$B\\
    \hline

    \theutterance \stepcounter{utterance}  
    & & \multicolumn{4}{p{0.6\linewidth}}{
        \cellcolor[rgb]{0.9,0.9,0.9}{
            \makecell[{{p{\linewidth}}}]{
                \texttt{\tiny{[P1$\langle$GM]}}
                \texttt{Jsem váš herní mistr a hrajete společnou hru s následující mřížkou jako herním plánem:} \\
\\ 
\texttt{```} \\
\texttt{ 1234567} \\
\texttt{╔════╤══╗} \\
\texttt{║◌◌◌O│PM║ 1} \\
\texttt{║◌◌T◌├──╢ 2} \\
\texttt{║◌◌◌◌│◌◌║ 3} \\
\texttt{║S┌──┼──╢ 4} \\
\texttt{║◌│◌A│◌◌║ 5} \\
\texttt{╟─┴┬─┤◌◌║ 6} \\
\texttt{║◌◌│◌│D◌║ 7} \\
\texttt{╚══╧═╧══╝} \\
\\ 
\texttt{```} \\
\\ 
\texttt{* Horní okraj zobrazuje souřadnice x rostoucí doprava a pravý okraj souřadnice y rostoucí dolů.} \\
\texttt{* Na vaší mřížce jsou náhodně umístěny následující objekty: 'P', 'O', 'T', 'S', 'D', 'A', 'M'.} \\
\\ 
\\ 
\texttt{Druhý hráč vidí variantu herní desky, kde jsou objekty umístěny na různých náhodných místech. Vy nevidíte desku druhého hráče a on nevidí tu vaši.} \\
\\ 
\\ 
\texttt{**Cíl:**} \\
\\ 
\texttt{Oba hráči musí přesunout objekty na svém pozadí tak, aby stejné objekty skončily na stejných souřadnicích. Musíte komunikovat s druhým hráčem, abyste se dohodli na společné cílové konfiguraci.} \\
\\ 
\\ 
\texttt{**Pravidla:**} \\
\\ 
\texttt{* V každém tahu můžete poslat přesně jeden z následujících dvou příkazů:} \\
\texttt{ 1. `ŘEKNI: <ZPRÁVA>`: odeslání zprávy (vše až do dalšího zalomení řádku) druhému hráči. Přepošlu ji vašemu partnerovi.} \\
\texttt{2. `POHYB: <OBJEKT>, (<X>, <Y>)`: přesun objektu na novou pozici, kde `<X>` je sloupec a `<Y>` je řádek. Budu vás informovat, zda byl váš přesun úspěšný, nebo ne.} \\
\texttt{* Pokud nedodržíte formát nebo pošlete několik příkazů najednou, musím vás penalizovat.} \\
\texttt{* Pokud oba hráči nasbírají více než 16 penalizací, oba prohrajete hru.} \\
\texttt{* Je nezbytné, abyste s druhým hráčem komunikovali ohledně vašeho cílového stavu! *Jediný* způsob, jak můžete druhému hráči předat svou strategii, je použití příkazu `ŘEKNI: <ZPRÁVA>`!} \\
\\ 
\\ 
\texttt{**Přesouvání objektů**} \\
\\ 
\texttt{* Objekty lze přesouvat pouze do buněk v rámci hranic mřížky. Cílová buňka musí být prázdná, tj. musí obsahovat pouze symbol '◌'.} \\
\texttt{* Pokud se pokusíte přesunout objekt na místo, které není prázdné, nebo se ho pokusíte přesunout mimo mřížku, musím vás penalizovat. Dostanete další pokus.} \\
\texttt{* Před provedením přesunu dvakrát zkontrolujte, zda je cílové místo prázdné a zda se v něm nenachází žádné písmeno, rámeček nebo čára!} \\
\\ 
\\ 
\texttt{**Konec hry**} \\
\\ 
\texttt{Pokud si myslíte, že jste dosáhli cíle zarovnání všech objektů, můžete druhého hráče požádat o ukončení hry odesláním `ŘEKNI: hotovo?`. Pokud vás druhý hráč požádá o ukončení hry a vy odpovíte `ŘEKNI: hotovo!`, hra skončí.} \\
\\ 
\texttt{Oba hráči vyhrají, pokud hra skončí během 28 kol, přičemž jedno kolo je definováno jako odeslání platného příkazu oběma hráči.} \\
\\ 
\texttt{**Bodování:**} \\
\\ 
\texttt{Čím blíže jsou stejné objekty na obou herních deskách, tím více bodů získáte. Penalizace vám body snižují. Podaří se vám překonat rekord?} \\
\\ 
\\ 
\texttt{Pošlete prosím zprávu druhému hráči a začněte hru!} \\
            }
        }
    }
    & \\ \\

    \theutterance \stepcounter{utterance}  
    & \multicolumn{4}{p{0.6\linewidth}}{
        \cellcolor[rgb]{0.9,0.9,0.9}{
            \makecell[{{p{\linewidth}}}]{
                \texttt{\tiny{[P1$\rangle$GM]}}
                \texttt{ŘEKNI: Ahoj! Přehrál jsem své objekty na mřížce. Zde jsou jejich počáteční souřadnice (X, Y): P: (6,1), O: (4,1), T: (3,2), S: (1,4), D: (6,7), A: (4,5), M: (7,1). Navrhuji dohodnout se na společném cílovém stavu. Můžeme zkusit seskupit je například do středu mřížky. Jaké jsou tvoje počáteční pozice?} \\
            }
        }
    }
    & & \\ \\

    \theutterance \stepcounter{utterance}  
    & & & \multicolumn{2}{p{0.3\linewidth}}{
        \cellcolor[rgb]{0.9,0.9,0.9}{
            \makecell[{{p{\linewidth}}}]{
                \texttt{\tiny{[GM$|$GM]}}
                \texttt{ŘEKNI: Ahoj! Přehrál jsem své objekty na mřížce. Zde jsou jejich počáteční souřadnice (X, Y): P: (6,1), O: (4,1), T: (3,2), S: (1,4), D: (6,7), A: (4,5), M: (7,1). Navrhuji dohodnout se na společném cílovém stavu. Můžeme zkusit seskupit je například do středu mřížky. Jaké jsou tvoje počáteční pozice?} \\
            }
        }
    }
    & & \\ \\

    \theutterance \stepcounter{utterance}  
    & & \multicolumn{4}{p{0.6\linewidth}}{
        \cellcolor[rgb]{0.9,0.9,0.9}{
            \makecell[{{p{\linewidth}}}]{
                \texttt{\tiny{[GM$\rangle$P2]}}
                \texttt{Jsem váš herní mistr a hrajete společnou hru s následující mřížkou jako herním plánem:} \\
\\ 
\texttt{```} \\
\texttt{ 1234567} \\
\texttt{╔════╤══╗} \\
\texttt{║◌◌◌◌│O◌║ 1} \\
\texttt{║P◌M◌├──╢ 2} \\
\texttt{║◌S◌◌│T◌║ 3} \\
\texttt{║◌┌──┼──╢ 4} \\
\texttt{║A│◌◌│◌◌║ 5} \\
\texttt{╟─┴┬─┤◌D║ 6} \\
\texttt{║◌◌│◌│◌◌║ 7} \\
\texttt{╚══╧═╧══╝} \\
\\ 
\texttt{```} \\
\\ 
\texttt{* Horní okraj zobrazuje souřadnice x rostoucí doprava a pravý okraj souřadnice y rostoucí dolů.} \\
\texttt{* Na vaší mřížce jsou náhodně umístěny následující objekty: 'P', 'O', 'T', 'S', 'D', 'A', 'M'.} \\
\\ 
\\ 
\texttt{Druhý hráč vidí variantu herní desky, kde jsou objekty umístěny na různých náhodných místech. Vy nevidíte desku druhého hráče a on nevidí tu vaši.} \\
\\ 
\\ 
\texttt{**Cíl:**} \\
\\ 
\texttt{Oba hráči musí přesunout objekty na svém pozadí tak, aby stejné objekty skončily na stejných souřadnicích. Musíte komunikovat s druhým hráčem, abyste se dohodli na společné cílové konfiguraci.} \\
\\ 
\\ 
\texttt{**Pravidla:**} \\
\\ 
\texttt{* V každém tahu můžete poslat přesně jeden z následujících dvou příkazů:} \\
\texttt{ 1. `ŘEKNI: <ZPRÁVA>`: odeslání zprávy (vše až do dalšího zalomení řádku) druhému hráči. Přepošlu ji vašemu partnerovi.} \\
\texttt{2. `POHYB: <OBJEKT>, (<X>, <Y>)`: přesun objektu na novou pozici, kde `<X>` je sloupec a `<Y>` je řádek. Budu vás informovat, zda byl váš přesun úspěšný, nebo ne.} \\
\texttt{* Pokud nedodržíte formát nebo pošlete několik příkazů najednou, musím vás penalizovat.} \\
\texttt{* Pokud oba hráči nasbírají více než 16 penalizací, oba prohrajete hru.} \\
\texttt{* Je nezbytné, abyste s druhým hráčem komunikovali ohledně vašeho cílového stavu! *Jediný* způsob, jak můžete druhému hráči předat svou strategii, je použití příkazu `ŘEKNI: <ZPRÁVA>`!} \\
\\ 
\\ 
\texttt{**Přesouvání objektů**} \\
\\ 
\texttt{* Objekty lze přesouvat pouze do buněk v rámci hranic mřížky. Cílová buňka musí být prázdná, tj. musí obsahovat pouze symbol '◌'.} \\
\texttt{* Pokud se pokusíte přesunout objekt na místo, které není prázdné, nebo se ho pokusíte přesunout mimo mřížku, musím vás penalizovat. Dostanete další pokus.} \\
\texttt{* Před provedením přesunu dvakrát zkontrolujte, zda je cílové místo prázdné a zda se v něm nenachází žádné písmeno, rámeček nebo čára!} \\
\\ 
\\ 
\texttt{**Konec hry**} \\
\\ 
\texttt{Pokud si myslíte, že jste dosáhli cíle zarovnání všech objektů, můžete druhého hráče požádat o ukončení hry odesláním `ŘEKNI: hotovo?`. Pokud vás druhý hráč požádá o ukončení hry a vy odpovíte `ŘEKNI: hotovo!`, hra skončí.} \\
\\ 
\texttt{Oba hráči vyhrají, pokud hra skončí během 28 kol, přičemž jedno kolo je definováno jako odeslání platného příkazu oběma hráči.} \\
\\ 
\texttt{**Bodování:**} \\
\\ 
\texttt{Čím blíže jsou stejné objekty na obou herních deskách, tím více bodů získáte. Penalizace vám body snižují. Podaří se vám překonat rekord?} \\
\\ 
\texttt{Druhý hráč zahájil hru odesláním této zprávy:} \\
\\ 
\texttt{"Ahoj! Přehrál jsem své objekty na mřížce. Zde jsou jejich počáteční souřadnice (X, Y): P: (6,1), O: (4,1), T: (3,2), S: (1,4), D: (6,7), A: (4,5), M: (7,1). Navrhuji dohodnout se na společném cílovém stavu. Můžeme zkusit seskupit je například do středu mřížky. Jaké jsou tvoje počáteční pozice?"} \\
\\ 
\texttt{Jaký je váš první příkaz?} \\
            }
        }
    }
    & \\ \\

    \theutterance \stepcounter{utterance}  
    & & & \multicolumn{4}{p{0.6\linewidth}}{
        \cellcolor[rgb]{0.9,0.9,0.9}{
            \makecell[{{p{\linewidth}}}]{
                \texttt{\tiny{[GM$\langle$P2]}}
                \texttt{ŘEKNI: Ahoj! Moje počáteční pozice (X, Y): P: (1,2), O: (5,1), T: (6,3), S: (3,3), D: (7,6), A: (1,5), M: (3,2). Souhlasím s dohodou na společném cílovém stavu. Vzhledem k překážkám na desce navrhuji jako společné cílové souřadnice následující volná pole: P: (1,1), O: (2,1), T: (3,1), S: (4,1), D: (6,1), A: (1,3), M: (2,2). Všechna tato pole vypadají jako prázdná bez čar. Jak se ti to zdí? Pokud souhlasíš, začnu přesouvat P na (1,1).} \\
            }
        }
    }
     \\ \\

    \theutterance \stepcounter{utterance}  
    & & & \multicolumn{2}{p{0.3\linewidth}}{
        \cellcolor[rgb]{0.9,0.9,0.9}{
            \makecell[{{p{\linewidth}}}]{
                \texttt{\tiny{[GM$|$GM]}}
                \texttt{ŘEKNI: Ahoj! Moje počáteční pozice (X, Y): P: (1,2), O: (5,1), T: (6,3), S: (3,3), D: (7,6), A: (1,5), M: (3,2). Souhlasím s dohodou na společném cílovém stavu. Vzhledem k překážkám na desce navrhuji jako společné cílové souřadnice následující volná pole: P: (1,1), O: (2,1), T: (3,1), S: (4,1), D: (6,1), A: (1,3), M: (2,2). Všechna tato pole vypadají jako prázdná bez čar. Jak se ti to zdí? Pokud souhlasíš, začnu přesouvat P na (1,1).} \\
            }
        }
    }
    & & \\ \\

    \theutterance \stepcounter{utterance}  
    & & \multicolumn{4}{p{0.6\linewidth}}{
        \cellcolor[rgb]{0.9,0.9,0.9}{
            \makecell[{{p{\linewidth}}}]{
                \texttt{\tiny{[P1$\langle$GM]}}
                \texttt{Vaše zpráva byla předána druhému hráči.} \\
\\ 
\\ 
\texttt{Právě hrajete kolo 1 z maximálně 28.} \\
\texttt{Dohromady jste nasbírali 0 z 16 pokut.} \\
\texttt{Druhý hráč poslal tuto zprávu:} \\
\\ 
\\ 
\texttt{"Ahoj! Moje počáteční pozice (X, Y): P: (1,2), O: (5,1), T: (6,3), S: (3,3), D: (7,6), A: (1,5), M: (3,2). Souhlasím s dohodou na společném cílovém stavu. Vzhledem k překážkám na desce navrhuji jako společné cílové souřadnice následující volná pole: P: (1,1), O: (2,1), T: (3,1), S: (4,1), D: (6,1), A: (1,3), M: (2,2). Všechna tato pole vypadají jako prázdná bez čar. Jak se ti to zdí? Pokud souhlasíš, začnu přesouvat P na (1,1)."} \\
\\ 
\\ 
\texttt{Jaký je váš další příkaz?} \\
            }
        }
    }
    & \\ \\

    \theutterance \stepcounter{utterance}  
    & \multicolumn{4}{p{0.6\linewidth}}{
        \cellcolor[rgb]{0.9,0.9,0.9}{
            \makecell[{{p{\linewidth}}}]{
                \texttt{\tiny{[P1$\rangle$GM]}}
                \texttt{ŘEKNI: Tvé cílové souřadnice nefungují, na mé desce jsou tam překážky (čáry/rámečky). Tady jsou volná prázdná pole (◌) na mé desce: (1,1), (1,2), (1,3), (2,2), (2,3), (3,1), (3,3), (5,3), (6,3), (7,3), (1,5), (3,5), (4,6), (5,6), (6,6), (7,6), (3,7), (5,7). Navrhuji jako společný cíl horní roh a spodní část: P:(1,1), O:(3,1), T:(1,2), S:(2,2), A:(1,3), D:(7,6), M:(3,3). Můžeš potvrdit, zda jsou na tvé desce tato pole volná?} \\
            }
        }
    }
    & & \\ \\

    \theutterance \stepcounter{utterance}  
    & & & \multicolumn{2}{p{0.3\linewidth}}{
        \cellcolor[rgb]{0.9,0.9,0.9}{
            \makecell[{{p{\linewidth}}}]{
                \texttt{\tiny{[GM$|$GM]}}
                \texttt{ŘEKNI: Tvé cílové souřadnice nefungují, na mé desce jsou tam překážky (čáry/rámečky). Tady jsou volná prázdná pole (◌) na mé desce: (1,1), (1,2), (1,3), (2,2), (2,3), (3,1), (3,3), (5,3), (6,3), (7,3), (1,5), (3,5), (4,6), (5,6), (6,6), (7,6), (3,7), (5,7). Navrhuji jako společný cíl horní roh a spodní část: P:(1,1), O:(3,1), T:(1,2), S:(2,2), A:(1,3), D:(7,6), M:(3,3). Můžeš potvrdit, zda jsou na tvé desce tato pole volná?} \\
            }
        }
    }
    & & \\ \\

    \theutterance \stepcounter{utterance}  
    & & \multicolumn{4}{p{0.6\linewidth}}{
        \cellcolor[rgb]{0.9,0.9,0.9}{
            \makecell[{{p{\linewidth}}}]{
                \texttt{\tiny{[GM$\rangle$P2]}}
                \texttt{Vaše zpráva byla předána druhému hráči.} \\
\\ 
\\ 
\texttt{Právě hrajete kolo 2 z maximálně 28.} \\
\texttt{Dohromady jste nasbírali 0 z 16 pokut.} \\
\texttt{Druhý hráč poslal tuto zprávu:} \\
\\ 
\\ 
\texttt{"Tvé cílové souřadnice nefungují, na mé desce jsou tam překážky (čáry/rámečky). Tady jsou volná prázdná pole (◌) na mé desce: (1,1), (1,2), (1,3), (2,2), (2,3), (3,1), (3,3), (5,3), (6,3), (7,3), (1,5), (3,5), (4,6), (5,6), (6,6), (7,6), (3,7), (5,7). Navrhuji jako společný cíl horní roh a spodní část: P:(1,1), O:(3,1), T:(1,2), S:(2,2), A:(1,3), D:(7,6), M:(3,3). Můžeš potvrdit, zda jsou na tvé desce tato pole volná?"} \\
\\ 
\\ 
\texttt{Jaký je váš další příkaz?} \\
            }
        }
    }
    & \\ \\

    \theutterance \stepcounter{utterance}  
    & & & \multicolumn{4}{p{0.6\linewidth}}{
        \cellcolor[rgb]{0.9,0.9,0.9}{
            \makecell[{{p{\linewidth}}}]{
                \texttt{\tiny{[GM$\langle$P2]}}
                \texttt{ŘEKNI: Zkontroloval jsem tvé navržené cílové souřadnice (P:(1,1), O:(3,1), T:(1,2), S:(2,2), A:(1,3), D:(7,6), M:(3,3)) a na mé desce jsou všechna tato pole volná (prázdné buňky ◌). Můj objekt M je sice aktuálně na (3,2), ale (3,3) je volná. Plně souhlasím s touto cílovou konfigurací. Teď přesunu svůj objekt P na pozici (1,1).} \\
            }
        }
    }
     \\ \\

    \theutterance \stepcounter{utterance}  
    & & & \multicolumn{2}{p{0.3\linewidth}}{
        \cellcolor[rgb]{0.9,0.9,0.9}{
            \makecell[{{p{\linewidth}}}]{
                \texttt{\tiny{[GM$|$GM]}}
                \texttt{ŘEKNI: Zkontroloval jsem tvé navržené cílové souřadnice (P:(1,1), O:(3,1), T:(1,2), S:(2,2), A:(1,3), D:(7,6), M:(3,3)) a na mé desce jsou všechna tato pole volná (prázdné buňky ◌). Můj objekt M je sice aktuálně na (3,2), ale (3,3) je volná. Plně souhlasím s touto cílovou konfigurací. Teď přesunu svůj objekt P na pozici (1,1).} \\
            }
        }
    }
    & & \\ \\

    \theutterance \stepcounter{utterance}  
    & & \multicolumn{4}{p{0.6\linewidth}}{
        \cellcolor[rgb]{0.9,0.9,0.9}{
            \makecell[{{p{\linewidth}}}]{
                \texttt{\tiny{[P1$\langle$GM]}}
                \texttt{Vaše zpráva byla předána druhému hráči.} \\
\\ 
\\ 
\texttt{Právě hrajete kolo 2 z maximálně 28.} \\
\texttt{Dohromady jste nasbírali 0 z 16 pokut.} \\
\texttt{Druhý hráč poslal tuto zprávu:} \\
\\ 
\\ 
\texttt{"Zkontroloval jsem tvé navržené cílové souřadnice (P:(1,1), O:(3,1), T:(1,2), S:(2,2), A:(1,3), D:(7,6), M:(3,3)) a na mé desce jsou všechna tato pole volná (prázdné buňky ◌). Můj objekt M je sice aktuálně na (3,2), ale (3,3) je volná. Plně souhlasím s touto cílovou konfigurací. Teď přesunu svůj objekt P na pozici (1,1)."} \\
\\ 
\\ 
\texttt{Jaký je váš další příkaz?} \\
            }
        }
    }
    & \\ \\

    \theutterance \stepcounter{utterance}  
    & \multicolumn{4}{p{0.6\linewidth}}{
        \cellcolor[rgb]{0.9,0.9,0.9}{
            \makecell[{{p{\linewidth}}}]{
                \texttt{\tiny{[P1$\rangle$GM]}}
                \texttt{POHYB: P, (1, 1)} \\
            }
        }
    }
    & & \\ \\

    \theutterance \stepcounter{utterance}  
    & & & \multicolumn{2}{p{0.3\linewidth}}{
        \cellcolor[rgb]{0.9,0.9,0.9}{
            \makecell[{{p{\linewidth}}}]{
                \texttt{\tiny{[GM$|$GM]}}
                \texttt{POHYB: P, (1, 1)} \\
            }
        }
    }
    & & \\ \\

    \theutterance \stepcounter{utterance}  
    & & & \multicolumn{2}{p{0.3\linewidth}}{
        \cellcolor[rgb]{0.9,0.9,0.9}{
            \makecell[{{p{\linewidth}}}]{
                \texttt{\tiny{[GM$|$GM]}}
                \texttt{P moved to (1, 1)} \\
            }
        }
    }
    & & \\ \\

    \theutterance \stepcounter{utterance}  
    & & \multicolumn{4}{p{0.6\linewidth}}{
        \cellcolor[rgb]{0.9,0.9,0.9}{
            \makecell[{{p{\linewidth}}}]{
                \texttt{\tiny{[GM$\rangle$P2]}}
                \texttt{Vaše zpráva byla předána druhému hráči.} \\
\\ 
\\ 
\texttt{Právě hrajete kolo 3 z maximálně 28.} \\
\texttt{Dohromady jste nasbírali 0 z 16 pokut.} \\
\texttt{Druhý hráč přesunul objekt na své mřížce.} \\
\\ 
\\ 
\texttt{Jaký je váš další příkaz?} \\
            }
        }
    }
    & \\ \\

    \theutterance \stepcounter{utterance}  
    & & & \multicolumn{4}{p{0.6\linewidth}}{
        \cellcolor[rgb]{0.9,0.9,0.9}{
            \makecell[{{p{\linewidth}}}]{
                \texttt{\tiny{[GM$\langle$P2]}}
                \texttt{POHYB: P, (1, 1)} \\
            }
        }
    }
     \\ \\

    \theutterance \stepcounter{utterance}  
    & & & \multicolumn{2}{p{0.3\linewidth}}{
        \cellcolor[rgb]{0.9,0.9,0.9}{
            \makecell[{{p{\linewidth}}}]{
                \texttt{\tiny{[GM$|$GM]}}
                \texttt{POHYB: P, (1, 1)} \\
            }
        }
    }
    & & \\ \\

    \theutterance \stepcounter{utterance}  
    & & & \multicolumn{2}{p{0.3\linewidth}}{
        \cellcolor[rgb]{0.9,0.9,0.9}{
            \makecell[{{p{\linewidth}}}]{
                \texttt{\tiny{[GM$|$GM]}}
                \texttt{P moved to (1, 1)} \\
            }
        }
    }
    & & \\ \\

    \theutterance \stepcounter{utterance}  
    & & \multicolumn{4}{p{0.6\linewidth}}{
        \cellcolor[rgb]{0.9,0.9,0.9}{
            \makecell[{{p{\linewidth}}}]{
                \texttt{\tiny{[P1$\langle$GM]}}
                \texttt{Úspěšně přesunuto 'P' do (1,1). Vaše aktualizovaná mřížka vypadá takto:} \\
\texttt{```} \\
\texttt{ 1234567} \\
\texttt{╔════╤══╗} \\
\texttt{║P◌◌O│◌M║ 1} \\
\texttt{║◌◌T◌├──╢ 2} \\
\texttt{║◌◌◌◌│◌◌║ 3} \\
\texttt{║S┌──┼──╢ 4} \\
\texttt{║◌│◌A│◌◌║ 5} \\
\texttt{╟─┴┬─┤◌◌║ 6} \\
\texttt{║◌◌│◌│D◌║ 7} \\
\texttt{╚══╧═╧══╝} \\
\texttt{```} \\
\texttt{Právě hrajete kolo 3 z maximálně 28.} \\
\texttt{Dohromady jste nasbírali 0 z 16 pokut.} \\
\texttt{Druhý hráč přesunul objekt na své mřížce.} \\
\\ 
\\ 
\texttt{Jaký je váš další příkaz?} \\
            }
        }
    }
    & \\ \\

    \theutterance \stepcounter{utterance}  
    & \multicolumn{4}{p{0.6\linewidth}}{
        \cellcolor[rgb]{0.9,0.9,0.9}{
            \makecell[{{p{\linewidth}}}]{
                \texttt{\tiny{[P1$\rangle$GM]}}
                \texttt{POHYB: O, (3, 1)} \\
            }
        }
    }
    & & \\ \\

    \theutterance \stepcounter{utterance}  
    & & & \multicolumn{2}{p{0.3\linewidth}}{
        \cellcolor[rgb]{0.9,0.9,0.9}{
            \makecell[{{p{\linewidth}}}]{
                \texttt{\tiny{[GM$|$GM]}}
                \texttt{POHYB: O, (3, 1)} \\
            }
        }
    }
    & & \\ \\

    \theutterance \stepcounter{utterance}  
    & & & \multicolumn{2}{p{0.3\linewidth}}{
        \cellcolor[rgb]{0.9,0.9,0.9}{
            \makecell[{{p{\linewidth}}}]{
                \texttt{\tiny{[GM$|$GM]}}
                \texttt{O moved to (3, 1)} \\
            }
        }
    }
    & & \\ \\

    \theutterance \stepcounter{utterance}  
    & & \multicolumn{4}{p{0.6\linewidth}}{
        \cellcolor[rgb]{0.9,0.9,0.9}{
            \makecell[{{p{\linewidth}}}]{
                \texttt{\tiny{[GM$\rangle$P2]}}
                \texttt{Úspěšně přesunuto 'P' do (1,1). Vaše aktualizovaná mřížka vypadá takto:} \\
\texttt{```} \\
\texttt{ 1234567} \\
\texttt{╔════╤══╗} \\
\texttt{║P◌◌◌│O◌║ 1} \\
\texttt{║◌◌M◌├──╢ 2} \\
\texttt{║◌S◌◌│T◌║ 3} \\
\texttt{║◌┌──┼──╢ 4} \\
\texttt{║A│◌◌│◌◌║ 5} \\
\texttt{╟─┴┬─┤◌D║ 6} \\
\texttt{║◌◌│◌│◌◌║ 7} \\
\texttt{╚══╧═╧══╝} \\
\texttt{```} \\
\texttt{Právě hrajete kolo 4 z maximálně 28.} \\
\texttt{Dohromady jste nasbírali 0 z 16 pokut.} \\
\texttt{Druhý hráč přesunul objekt na své mřížce.} \\
\\ 
\\ 
\texttt{Jaký je váš další příkaz?} \\
            }
        }
    }
    & \\ \\

    \theutterance \stepcounter{utterance}  
    & & & \multicolumn{4}{p{0.6\linewidth}}{
        \cellcolor[rgb]{0.9,0.9,0.9}{
            \makecell[{{p{\linewidth}}}]{
                \texttt{\tiny{[GM$\langle$P2]}}
                \texttt{POHYB: O, (3, 1)} \\
            }
        }
    }
     \\ \\

    \theutterance \stepcounter{utterance}  
    & & & \multicolumn{2}{p{0.3\linewidth}}{
        \cellcolor[rgb]{0.9,0.9,0.9}{
            \makecell[{{p{\linewidth}}}]{
                \texttt{\tiny{[GM$|$GM]}}
                \texttt{POHYB: O, (3, 1)} \\
            }
        }
    }
    & & \\ \\

    \theutterance \stepcounter{utterance}  
    & & & \multicolumn{2}{p{0.3\linewidth}}{
        \cellcolor[rgb]{0.9,0.9,0.9}{
            \makecell[{{p{\linewidth}}}]{
                \texttt{\tiny{[GM$|$GM]}}
                \texttt{O moved to (3, 1)} \\
            }
        }
    }
    & & \\ \\

    \theutterance \stepcounter{utterance}  
    & & \multicolumn{4}{p{0.6\linewidth}}{
        \cellcolor[rgb]{0.9,0.9,0.9}{
            \makecell[{{p{\linewidth}}}]{
                \texttt{\tiny{[P1$\langle$GM]}}
                \texttt{Úspěšně přesunuto 'O' do (3,1). Vaše aktualizovaná mřížka vypadá takto:} \\
\texttt{```} \\
\texttt{ 1234567} \\
\texttt{╔════╤══╗} \\
\texttt{║P◌O◌│◌M║ 1} \\
\texttt{║◌◌T◌├──╢ 2} \\
\texttt{║◌◌◌◌│◌◌║ 3} \\
\texttt{║S┌──┼──╢ 4} \\
\texttt{║◌│◌A│◌◌║ 5} \\
\texttt{╟─┴┬─┤◌◌║ 6} \\
\texttt{║◌◌│◌│D◌║ 7} \\
\texttt{╚══╧═╧══╝} \\
\texttt{```} \\
\texttt{Právě hrajete kolo 4 z maximálně 28.} \\
\texttt{Dohromady jste nasbírali 0 z 16 pokut.} \\
\texttt{Druhý hráč přesunul objekt na své mřížce.} \\
\\ 
\\ 
\texttt{Jaký je váš další příkaz?} \\
            }
        }
    }
    & \\ \\

    \theutterance \stepcounter{utterance}  
    & \multicolumn{4}{p{0.6\linewidth}}{
        \cellcolor[rgb]{0.9,0.9,0.9}{
            \makecell[{{p{\linewidth}}}]{
                \texttt{\tiny{[P1$\rangle$GM]}}
                \texttt{POHYB: T, (1, 2)} \\
            }
        }
    }
    & & \\ \\

    \theutterance \stepcounter{utterance}  
    & & & \multicolumn{2}{p{0.3\linewidth}}{
        \cellcolor[rgb]{0.9,0.9,0.9}{
            \makecell[{{p{\linewidth}}}]{
                \texttt{\tiny{[GM$|$GM]}}
                \texttt{POHYB: T, (1, 2)} \\
            }
        }
    }
    & & \\ \\

    \theutterance \stepcounter{utterance}  
    & & & \multicolumn{2}{p{0.3\linewidth}}{
        \cellcolor[rgb]{0.9,0.9,0.9}{
            \makecell[{{p{\linewidth}}}]{
                \texttt{\tiny{[GM$|$GM]}}
                \texttt{T moved to (1, 2)} \\
            }
        }
    }
    & & \\ \\

    \theutterance \stepcounter{utterance}  
    & & \multicolumn{4}{p{0.6\linewidth}}{
        \cellcolor[rgb]{0.9,0.9,0.9}{
            \makecell[{{p{\linewidth}}}]{
                \texttt{\tiny{[GM$\rangle$P2]}}
                \texttt{Úspěšně přesunuto 'O' do (3,1). Vaše aktualizovaná mřížka vypadá takto:} \\
\texttt{```} \\
\texttt{ 1234567} \\
\texttt{╔════╤══╗} \\
\texttt{║P◌O◌│◌◌║ 1} \\
\texttt{║◌◌M◌├──╢ 2} \\
\texttt{║◌S◌◌│T◌║ 3} \\
\texttt{║◌┌──┼──╢ 4} \\
\texttt{║A│◌◌│◌◌║ 5} \\
\texttt{╟─┴┬─┤◌D║ 6} \\
\texttt{║◌◌│◌│◌◌║ 7} \\
\texttt{╚══╧═╧══╝} \\
\texttt{```} \\
\texttt{Právě hrajete kolo 5 z maximálně 28.} \\
\texttt{Dohromady jste nasbírali 0 z 16 pokut.} \\
\texttt{Druhý hráč přesunul objekt na své mřížce.} \\
\\ 
\\ 
\texttt{Jaký je váš další příkaz?} \\
            }
        }
    }
    & \\ \\

    \theutterance \stepcounter{utterance}  
    & & & \multicolumn{4}{p{0.6\linewidth}}{
        \cellcolor[rgb]{0.9,0.9,0.9}{
            \makecell[{{p{\linewidth}}}]{
                \texttt{\tiny{[GM$\langle$P2]}}
                \texttt{POHYB: T, (1, 2)} \\
            }
        }
    }
     \\ \\

    \theutterance \stepcounter{utterance}  
    & & & \multicolumn{2}{p{0.3\linewidth}}{
        \cellcolor[rgb]{0.9,0.9,0.9}{
            \makecell[{{p{\linewidth}}}]{
                \texttt{\tiny{[GM$|$GM]}}
                \texttt{POHYB: T, (1, 2)} \\
            }
        }
    }
    & & \\ \\

    \theutterance \stepcounter{utterance}  
    & & & \multicolumn{2}{p{0.3\linewidth}}{
        \cellcolor[rgb]{0.9,0.9,0.9}{
            \makecell[{{p{\linewidth}}}]{
                \texttt{\tiny{[GM$|$GM]}}
                \texttt{T moved to (1, 2)} \\
            }
        }
    }
    & & \\ \\

    \theutterance \stepcounter{utterance}  
    & & \multicolumn{4}{p{0.6\linewidth}}{
        \cellcolor[rgb]{0.9,0.9,0.9}{
            \makecell[{{p{\linewidth}}}]{
                \texttt{\tiny{[P1$\langle$GM]}}
                \texttt{Úspěšně přesunuto 'T' do (1,2). Vaše aktualizovaná mřížka vypadá takto:} \\
\texttt{```} \\
\texttt{ 1234567} \\
\texttt{╔════╤══╗} \\
\texttt{║P◌O◌│◌M║ 1} \\
\texttt{║T◌◌◌├──╢ 2} \\
\texttt{║◌◌◌◌│◌◌║ 3} \\
\texttt{║S┌──┼──╢ 4} \\
\texttt{║◌│◌A│◌◌║ 5} \\
\texttt{╟─┴┬─┤◌◌║ 6} \\
\texttt{║◌◌│◌│D◌║ 7} \\
\texttt{╚══╧═╧══╝} \\
\texttt{```} \\
\texttt{Právě hrajete kolo 5 z maximálně 28.} \\
\texttt{Dohromady jste nasbírali 0 z 16 pokut.} \\
\texttt{Druhý hráč přesunul objekt na své mřížce.} \\
\\ 
\\ 
\texttt{Jaký je váš další příkaz?} \\
            }
        }
    }
    & \\ \\

    \theutterance \stepcounter{utterance}  
    & \multicolumn{4}{p{0.6\linewidth}}{
        \cellcolor[rgb]{0.9,0.9,0.9}{
            \makecell[{{p{\linewidth}}}]{
                \texttt{\tiny{[P1$\rangle$GM]}}
                \texttt{POHYB: S, (2, 2)} \\
            }
        }
    }
    & & \\ \\

    \theutterance \stepcounter{utterance}  
    & & & \multicolumn{2}{p{0.3\linewidth}}{
        \cellcolor[rgb]{0.9,0.9,0.9}{
            \makecell[{{p{\linewidth}}}]{
                \texttt{\tiny{[GM$|$GM]}}
                \texttt{POHYB: S, (2, 2)} \\
            }
        }
    }
    & & \\ \\

    \theutterance \stepcounter{utterance}  
    & & & \multicolumn{2}{p{0.3\linewidth}}{
        \cellcolor[rgb]{0.9,0.9,0.9}{
            \makecell[{{p{\linewidth}}}]{
                \texttt{\tiny{[GM$|$GM]}}
                \texttt{S moved to (2, 2)} \\
            }
        }
    }
    & & \\ \\

    \theutterance \stepcounter{utterance}  
    & & \multicolumn{4}{p{0.6\linewidth}}{
        \cellcolor[rgb]{0.9,0.9,0.9}{
            \makecell[{{p{\linewidth}}}]{
                \texttt{\tiny{[GM$\rangle$P2]}}
                \texttt{Úspěšně přesunuto 'T' do (1,2). Vaše aktualizovaná mřížka vypadá takto:} \\
\texttt{```} \\
\texttt{ 1234567} \\
\texttt{╔════╤══╗} \\
\texttt{║P◌O◌│◌◌║ 1} \\
\texttt{║T◌M◌├──╢ 2} \\
\texttt{║◌S◌◌│◌◌║ 3} \\
\texttt{║◌┌──┼──╢ 4} \\
\texttt{║A│◌◌│◌◌║ 5} \\
\texttt{╟─┴┬─┤◌D║ 6} \\
\texttt{║◌◌│◌│◌◌║ 7} \\
\texttt{╚══╧═╧══╝} \\
\texttt{```} \\
\texttt{Právě hrajete kolo 6 z maximálně 28.} \\
\texttt{Dohromady jste nasbírali 0 z 16 pokut.} \\
\texttt{Druhý hráč přesunul objekt na své mřížce.} \\
\\ 
\\ 
\texttt{Jaký je váš další příkaz?} \\
            }
        }
    }
    & \\ \\

    \theutterance \stepcounter{utterance}  
    & & & \multicolumn{4}{p{0.6\linewidth}}{
        \cellcolor[rgb]{0.9,0.9,0.9}{
            \makecell[{{p{\linewidth}}}]{
                \texttt{\tiny{[GM$\langle$P2]}}
                \texttt{POHYB: S, (2, 2)} \\
            }
        }
    }
     \\ \\

    \theutterance \stepcounter{utterance}  
    & & & \multicolumn{2}{p{0.3\linewidth}}{
        \cellcolor[rgb]{0.9,0.9,0.9}{
            \makecell[{{p{\linewidth}}}]{
                \texttt{\tiny{[GM$|$GM]}}
                \texttt{POHYB: S, (2, 2)} \\
            }
        }
    }
    & & \\ \\

    \theutterance \stepcounter{utterance}  
    & & & \multicolumn{2}{p{0.3\linewidth}}{
        \cellcolor[rgb]{0.9,0.9,0.9}{
            \makecell[{{p{\linewidth}}}]{
                \texttt{\tiny{[GM$|$GM]}}
                \texttt{S moved to (2, 2)} \\
            }
        }
    }
    & & \\ \\

    \theutterance \stepcounter{utterance}  
    & & \multicolumn{4}{p{0.6\linewidth}}{
        \cellcolor[rgb]{0.9,0.9,0.9}{
            \makecell[{{p{\linewidth}}}]{
                \texttt{\tiny{[P1$\langle$GM]}}
                \texttt{Úspěšně přesunuto 'S' do (2,2). Vaše aktualizovaná mřížka vypadá takto:} \\
\texttt{```} \\
\texttt{ 1234567} \\
\texttt{╔════╤══╗} \\
\texttt{║P◌O◌│◌M║ 1} \\
\texttt{║TS◌◌├──╢ 2} \\
\texttt{║◌◌◌◌│◌◌║ 3} \\
\texttt{║◌┌──┼──╢ 4} \\
\texttt{║◌│◌A│◌◌║ 5} \\
\texttt{╟─┴┬─┤◌◌║ 6} \\
\texttt{║◌◌│◌│D◌║ 7} \\
\texttt{╚══╧═╧══╝} \\
\texttt{```} \\
\texttt{Právě hrajete kolo 6 z maximálně 28.} \\
\texttt{Dohromady jste nasbírali 0 z 16 pokut.} \\
\texttt{Druhý hráč přesunul objekt na své mřížce.} \\
\\ 
\\ 
\texttt{Jaký je váš další příkaz?} \\
            }
        }
    }
    & \\ \\

    \theutterance \stepcounter{utterance}  
    & \multicolumn{4}{p{0.6\linewidth}}{
        \cellcolor[rgb]{0.9,0.9,0.9}{
            \makecell[{{p{\linewidth}}}]{
                \texttt{\tiny{[P1$\rangle$GM]}}
                \texttt{POHYB: A, (1, 3)} \\
            }
        }
    }
    & & \\ \\

    \theutterance \stepcounter{utterance}  
    & & & \multicolumn{2}{p{0.3\linewidth}}{
        \cellcolor[rgb]{0.9,0.9,0.9}{
            \makecell[{{p{\linewidth}}}]{
                \texttt{\tiny{[GM$|$GM]}}
                \texttt{POHYB: A, (1, 3)} \\
            }
        }
    }
    & & \\ \\

    \theutterance \stepcounter{utterance}  
    & & & \multicolumn{2}{p{0.3\linewidth}}{
        \cellcolor[rgb]{0.9,0.9,0.9}{
            \makecell[{{p{\linewidth}}}]{
                \texttt{\tiny{[GM$|$GM]}}
                \texttt{A moved to (1, 3)} \\
            }
        }
    }
    & & \\ \\

    \theutterance \stepcounter{utterance}  
    & & \multicolumn{4}{p{0.6\linewidth}}{
        \cellcolor[rgb]{0.9,0.9,0.9}{
            \makecell[{{p{\linewidth}}}]{
                \texttt{\tiny{[GM$\rangle$P2]}}
                \texttt{Úspěšně přesunuto 'S' do (2,2). Vaše aktualizovaná mřížka vypadá takto:} \\
\texttt{```} \\
\texttt{ 1234567} \\
\texttt{╔════╤══╗} \\
\texttt{║P◌O◌│◌◌║ 1} \\
\texttt{║TSM◌├──╢ 2} \\
\texttt{║◌◌◌◌│◌◌║ 3} \\
\texttt{║◌┌──┼──╢ 4} \\
\texttt{║A│◌◌│◌◌║ 5} \\
\texttt{╟─┴┬─┤◌D║ 6} \\
\texttt{║◌◌│◌│◌◌║ 7} \\
\texttt{╚══╧═╧══╝} \\
\texttt{```} \\
\texttt{Právě hrajete kolo 7 z maximálně 28.} \\
\texttt{Dohromady jste nasbírali 0 z 16 pokut.} \\
\texttt{Druhý hráč přesunul objekt na své mřížce.} \\
\\ 
\\ 
\texttt{Jaký je váš další příkaz?} \\
            }
        }
    }
    & \\ \\

    \theutterance \stepcounter{utterance}  
    & & & \multicolumn{4}{p{0.6\linewidth}}{
        \cellcolor[rgb]{0.9,0.9,0.9}{
            \makecell[{{p{\linewidth}}}]{
                \texttt{\tiny{[GM$\langle$P2]}}
                \texttt{POHYB: A, (1, 3)} \\
            }
        }
    }
     \\ \\

    \theutterance \stepcounter{utterance}  
    & & & \multicolumn{2}{p{0.3\linewidth}}{
        \cellcolor[rgb]{0.9,0.9,0.9}{
            \makecell[{{p{\linewidth}}}]{
                \texttt{\tiny{[GM$|$GM]}}
                \texttt{POHYB: A, (1, 3)} \\
            }
        }
    }
    & & \\ \\

    \theutterance \stepcounter{utterance}  
    & & & \multicolumn{2}{p{0.3\linewidth}}{
        \cellcolor[rgb]{0.9,0.9,0.9}{
            \makecell[{{p{\linewidth}}}]{
                \texttt{\tiny{[GM$|$GM]}}
                \texttt{A moved to (1, 3)} \\
            }
        }
    }
    & & \\ \\

    \theutterance \stepcounter{utterance}  
    & & \multicolumn{4}{p{0.6\linewidth}}{
        \cellcolor[rgb]{0.9,0.9,0.9}{
            \makecell[{{p{\linewidth}}}]{
                \texttt{\tiny{[P1$\langle$GM]}}
                \texttt{Úspěšně přesunuto 'A' do (1,3). Vaše aktualizovaná mřížka vypadá takto:} \\
\texttt{```} \\
\texttt{ 1234567} \\
\texttt{╔════╤══╗} \\
\texttt{║P◌O◌│◌M║ 1} \\
\texttt{║TS◌◌├──╢ 2} \\
\texttt{║A◌◌◌│◌◌║ 3} \\
\texttt{║◌┌──┼──╢ 4} \\
\texttt{║◌│◌◌│◌◌║ 5} \\
\texttt{╟─┴┬─┤◌◌║ 6} \\
\texttt{║◌◌│◌│D◌║ 7} \\
\texttt{╚══╧═╧══╝} \\
\texttt{```} \\
\texttt{Právě hrajete kolo 7 z maximálně 28.} \\
\texttt{Dohromady jste nasbírali 0 z 16 pokut.} \\
\texttt{Druhý hráč přesunul objekt na své mřížce.} \\
\\ 
\\ 
\texttt{Jaký je váš další příkaz?} \\
            }
        }
    }
    & \\ \\

    \theutterance \stepcounter{utterance}  
    & \multicolumn{4}{p{0.6\linewidth}}{
        \cellcolor[rgb]{0.9,0.9,0.9}{
            \makecell[{{p{\linewidth}}}]{
                \texttt{\tiny{[P1$\rangle$GM]}}
                \texttt{POHYB: D, (7, 6)} \\
            }
        }
    }
    & & \\ \\

    \theutterance \stepcounter{utterance}  
    & & & \multicolumn{2}{p{0.3\linewidth}}{
        \cellcolor[rgb]{0.9,0.9,0.9}{
            \makecell[{{p{\linewidth}}}]{
                \texttt{\tiny{[GM$|$GM]}}
                \texttt{POHYB: D, (7, 6)} \\
            }
        }
    }
    & & \\ \\

    \theutterance \stepcounter{utterance}  
    & & & \multicolumn{2}{p{0.3\linewidth}}{
        \cellcolor[rgb]{0.9,0.9,0.9}{
            \makecell[{{p{\linewidth}}}]{
                \texttt{\tiny{[GM$|$GM]}}
                \texttt{D moved to (7, 6)} \\
            }
        }
    }
    & & \\ \\

    \theutterance \stepcounter{utterance}  
    & & \multicolumn{4}{p{0.6\linewidth}}{
        \cellcolor[rgb]{0.9,0.9,0.9}{
            \makecell[{{p{\linewidth}}}]{
                \texttt{\tiny{[GM$\rangle$P2]}}
                \texttt{Úspěšně přesunuto 'A' do (1,3). Vaše aktualizovaná mřížka vypadá takto:} \\
\texttt{```} \\
\texttt{ 1234567} \\
\texttt{╔════╤══╗} \\
\texttt{║P◌O◌│◌◌║ 1} \\
\texttt{║TSM◌├──╢ 2} \\
\texttt{║A◌◌◌│◌◌║ 3} \\
\texttt{║◌┌──┼──╢ 4} \\
\texttt{║◌│◌◌│◌◌║ 5} \\
\texttt{╟─┴┬─┤◌D║ 6} \\
\texttt{║◌◌│◌│◌◌║ 7} \\
\texttt{╚══╧═╧══╝} \\
\texttt{```} \\
\texttt{Právě hrajete kolo 8 z maximálně 28.} \\
\texttt{Dohromady jste nasbírali 0 z 16 pokut.} \\
\texttt{Druhý hráč přesunul objekt na své mřížce.} \\
\\ 
\\ 
\texttt{Jaký je váš další příkaz?} \\
            }
        }
    }
    & \\ \\

    \theutterance \stepcounter{utterance}  
    & & & \multicolumn{4}{p{0.6\linewidth}}{
        \cellcolor[rgb]{0.9,0.9,0.9}{
            \makecell[{{p{\linewidth}}}]{
                \texttt{\tiny{[GM$\langle$P2]}}
                \texttt{POHYB: M, (3, 3)} \\
            }
        }
    }
     \\ \\

    \theutterance \stepcounter{utterance}  
    & & & \multicolumn{2}{p{0.3\linewidth}}{
        \cellcolor[rgb]{0.9,0.9,0.9}{
            \makecell[{{p{\linewidth}}}]{
                \texttt{\tiny{[GM$|$GM]}}
                \texttt{POHYB: M, (3, 3)} \\
            }
        }
    }
    & & \\ \\

    \theutterance \stepcounter{utterance}  
    & & & \multicolumn{2}{p{0.3\linewidth}}{
        \cellcolor[rgb]{0.9,0.9,0.9}{
            \makecell[{{p{\linewidth}}}]{
                \texttt{\tiny{[GM$|$GM]}}
                \texttt{M moved to (3, 3)} \\
            }
        }
    }
    & & \\ \\

    \theutterance \stepcounter{utterance}  
    & & \multicolumn{4}{p{0.6\linewidth}}{
        \cellcolor[rgb]{0.9,0.9,0.9}{
            \makecell[{{p{\linewidth}}}]{
                \texttt{\tiny{[P1$\langle$GM]}}
                \texttt{Úspěšně přesunuto 'D' do (7,6). Vaše aktualizovaná mřížka vypadá takto:} \\
\texttt{```} \\
\texttt{ 1234567} \\
\texttt{╔════╤══╗} \\
\texttt{║P◌O◌│◌M║ 1} \\
\texttt{║TS◌◌├──╢ 2} \\
\texttt{║A◌◌◌│◌◌║ 3} \\
\texttt{║◌┌──┼──╢ 4} \\
\texttt{║◌│◌◌│◌◌║ 5} \\
\texttt{╟─┴┬─┤◌D║ 6} \\
\texttt{║◌◌│◌│◌◌║ 7} \\
\texttt{╚══╧═╧══╝} \\
\texttt{```} \\
\texttt{Právě hrajete kolo 8 z maximálně 28.} \\
\texttt{Dohromady jste nasbírali 0 z 16 pokut.} \\
\texttt{Druhý hráč přesunul objekt na své mřížce.} \\
\\ 
\\ 
\texttt{Jaký je váš další příkaz?} \\
            }
        }
    }
    & \\ \\

    \theutterance \stepcounter{utterance}  
    & \multicolumn{4}{p{0.6\linewidth}}{
        \cellcolor[rgb]{0.9,0.9,0.9}{
            \makecell[{{p{\linewidth}}}]{
                \texttt{\tiny{[P1$\rangle$GM]}}
                \texttt{POHYB: M, (3, 3)} \\
            }
        }
    }
    & & \\ \\

    \theutterance \stepcounter{utterance}  
    & & & \multicolumn{2}{p{0.3\linewidth}}{
        \cellcolor[rgb]{0.9,0.9,0.9}{
            \makecell[{{p{\linewidth}}}]{
                \texttt{\tiny{[GM$|$GM]}}
                \texttt{POHYB: M, (3, 3)} \\
            }
        }
    }
    & & \\ \\

    \theutterance \stepcounter{utterance}  
    & & & \multicolumn{2}{p{0.3\linewidth}}{
        \cellcolor[rgb]{0.9,0.9,0.9}{
            \makecell[{{p{\linewidth}}}]{
                \texttt{\tiny{[GM$|$GM]}}
                \texttt{M moved to (3, 3)} \\
            }
        }
    }
    & & \\ \\

    \theutterance \stepcounter{utterance}  
    & & \multicolumn{4}{p{0.6\linewidth}}{
        \cellcolor[rgb]{0.9,0.9,0.9}{
            \makecell[{{p{\linewidth}}}]{
                \texttt{\tiny{[GM$\rangle$P2]}}
                \texttt{Úspěšně přesunuto 'M' do (3,3). Vaše aktualizovaná mřížka vypadá takto:} \\
\texttt{```} \\
\texttt{ 1234567} \\
\texttt{╔════╤══╗} \\
\texttt{║P◌O◌│◌◌║ 1} \\
\texttt{║TS◌◌├──╢ 2} \\
\texttt{║A◌M◌│◌◌║ 3} \\
\texttt{║◌┌──┼──╢ 4} \\
\texttt{║◌│◌◌│◌◌║ 5} \\
\texttt{╟─┴┬─┤◌D║ 6} \\
\texttt{║◌◌│◌│◌◌║ 7} \\
\texttt{╚══╧═╧══╝} \\
\texttt{```} \\
\texttt{Právě hrajete kolo 9 z maximálně 28.} \\
\texttt{Dohromady jste nasbírali 0 z 16 pokut.} \\
\texttt{Druhý hráč přesunul objekt na své mřížce.} \\
\\ 
\\ 
\texttt{Jaký je váš další příkaz?} \\
            }
        }
    }
    & \\ \\

    \theutterance \stepcounter{utterance}  
    & & & \multicolumn{4}{p{0.6\linewidth}}{
        \cellcolor[rgb]{0.9,0.9,0.9}{
            \makecell[{{p{\linewidth}}}]{
                \texttt{\tiny{[GM$\langle$P2]}}
                \texttt{POHYB: D, (7, 6)} \\
            }
        }
    }
     \\ \\

    \theutterance \stepcounter{utterance}  
    & & & \multicolumn{2}{p{0.3\linewidth}}{
        \cellcolor[rgb]{0.9,0.9,0.9}{
            \makecell[{{p{\linewidth}}}]{
                \texttt{\tiny{[GM$|$GM]}}
                \texttt{POHYB: D, (7, 6)} \\
            }
        }
    }
    & & \\ \\

    \theutterance \stepcounter{utterance}  
    & & & \multicolumn{2}{p{0.3\linewidth}}{
        \cellcolor[rgb]{0.9,0.9,0.9}{
            \makecell[{{p{\linewidth}}}]{
                \texttt{\tiny{[GM$|$GM]}}
                \texttt{Penalizace: (7,6) není prázdný, ale obsahuje 'D'.} \\
\\ 
\texttt{Dohromady jste nasbírali 1 z 16 pokut.} \\
\texttt{Zkuste to prosím znovu!} \\
            }
        }
    }
    & & \\ \\

    \theutterance \stepcounter{utterance}  
    & & \multicolumn{4}{p{0.6\linewidth}}{
        \cellcolor[rgb]{0.9,0.9,0.9}{
            \makecell[{{p{\linewidth}}}]{
                \texttt{\tiny{[GM$\rangle$P2]}}
                \texttt{Penalizace: (7,6) není prázdný, ale obsahuje 'D'.} \\
\\ 
\texttt{Dohromady jste nasbírali 1 z 16 pokut.} \\
\texttt{Zkuste to prosím znovu!} \\
            }
        }
    }
    & \\ \\

    \theutterance \stepcounter{utterance}  
    & & & \multicolumn{4}{p{0.6\linewidth}}{
        \cellcolor[rgb]{0.9,0.9,0.9}{
            \makecell[{{p{\linewidth}}}]{
                \texttt{\tiny{[GM$\langle$P2]}}
                \texttt{Ah, D už je na (7,6)! Omlouvám se za přehlédnutí. Všechny objekty jsou nyní na cílových pozicích. Pošlu zprávu druhému hráči, abychom zkontrolovali stav.} \\
\\ 
\texttt{ŘEKNI: Všechny mé objekty jsou nyní na dohodnutých cílových pozicích: P:(1,1), O:(3,1), T:(1,2), S:(2,2), A:(1,3), D:(7,6), M:(3,3). Jaký je tvůj stav? Jsou i tvé objekty na místech? Můžeme ukončit hru?} \\
            }
        }
    }
     \\ \\

    \theutterance \stepcounter{utterance}  
    & & & \multicolumn{2}{p{0.3\linewidth}}{
        \cellcolor[rgb]{0.9,0.9,0.9}{
            \makecell[{{p{\linewidth}}}]{
                \texttt{\tiny{[GM$|$GM]}}
                \texttt{Ah, D už je na (7,6)! Omlouvám se za přehlédnutí. Všechny objekty jsou nyní na cílových pozicích. Pošlu zprávu druhému hráči, abychom zkontrolovali stav.} \\
\\ 
\texttt{ŘEKNI: Všechny mé objekty jsou nyní na dohodnutých cílových pozicích: P:(1,1), O:(3,1), T:(1,2), S:(2,2), A:(1,3), D:(7,6), M:(3,3). Jaký je tvůj stav? Jsou i tvé objekty na místech? Můžeme ukončit hru?} \\
            }
        }
    }
    & & \\ \\

    \theutterance \stepcounter{utterance}  
    & & & \multicolumn{2}{p{0.3\linewidth}}{
        \cellcolor[rgb]{0.9,0.9,0.9}{
            \makecell[{{p{\linewidth}}}]{
                \texttt{\tiny{[GM$|$GM]}}
                \texttt{Invalid format: head is not empty: 'Ah, D už je na (7,6)! Omlouvám se za přehlédnutí. Všechny objekty jsou nyní na cílových pozicích. Pošlu zprávu druhému hráči, abychom zkontrolovali stav.} \\
\\ 
\texttt{'} \\
            }
        }
    }
    & & \\ \\

    \theutterance \stepcounter{utterance}  
    & & \multicolumn{4}{p{0.6\linewidth}}{
        \cellcolor[rgb]{0.9,0.9,0.9}{
            \makecell[{{p{\linewidth}}}]{
                \texttt{\tiny{[GM$\rangle$P2]}}
                \texttt{Trest: Vaše zpráva nesmí obsahovat nic před příkazem!} \\
\\ 
\texttt{Ujistěte se, že vaše odpověď obsahuje pouze `ŘEKNI: <ZPRÁVA>` nebo `POHYB: <OBJEKT>, (<X>, <Y>)` a nic jiného!} \\
\\ 
\texttt{Dohromady jste nasbírali 2 z 16 pokut.} \\
\texttt{Zkuste to prosím znovu!} \\
            }
        }
    }
    & \\ \\

    \theutterance \stepcounter{utterance}  
    & & & \multicolumn{4}{p{0.6\linewidth}}{
        \cellcolor[rgb]{0.9,0.9,0.9}{
            \makecell[{{p{\linewidth}}}]{
                \texttt{\tiny{[GM$\langle$P2]}}
                \texttt{ŘEKNI: Všechny mé objekty jsou nyní na dohodnutých cílových pozicích: P:(1,1), O:(3,1), T:(1,2), S:(2,2), A:(1,3), D:(7,6), M:(3,3). Jaký je tvůj stav? Jsou i tvé objekty na místech? Můžeme ukončit hru?} \\
            }
        }
    }
     \\ \\

    \theutterance \stepcounter{utterance}  
    & & & \multicolumn{2}{p{0.3\linewidth}}{
        \cellcolor[rgb]{0.9,0.9,0.9}{
            \makecell[{{p{\linewidth}}}]{
                \texttt{\tiny{[GM$|$GM]}}
                \texttt{ŘEKNI: Všechny mé objekty jsou nyní na dohodnutých cílových pozicích: P:(1,1), O:(3,1), T:(1,2), S:(2,2), A:(1,3), D:(7,6), M:(3,3). Jaký je tvůj stav? Jsou i tvé objekty na místech? Můžeme ukončit hru?} \\
            }
        }
    }
    & & \\ \\

    \theutterance \stepcounter{utterance}  
    & & \multicolumn{4}{p{0.6\linewidth}}{
        \cellcolor[rgb]{0.9,0.9,0.9}{
            \makecell[{{p{\linewidth}}}]{
                \texttt{\tiny{[P1$\langle$GM]}}
                \texttt{Úspěšně přesunuto 'M' do (3,3). Vaše aktualizovaná mřížka vypadá takto:} \\
\texttt{```} \\
\texttt{ 1234567} \\
\texttt{╔════╤══╗} \\
\texttt{║P◌O◌│◌◌║ 1} \\
\texttt{║TS◌◌├──╢ 2} \\
\texttt{║A◌M◌│◌◌║ 3} \\
\texttt{║◌┌──┼──╢ 4} \\
\texttt{║◌│◌◌│◌◌║ 5} \\
\texttt{╟─┴┬─┤◌D║ 6} \\
\texttt{║◌◌│◌│◌◌║ 7} \\
\texttt{╚══╧═╧══╝} \\
\texttt{```} \\
\texttt{Právě hrajete kolo 9 z maximálně 28.} \\
\texttt{Dohromady jste nasbírali 2 z 16 pokut.} \\
\texttt{Druhý hráč poslal tuto zprávu:} \\
\\ 
\\ 
\texttt{"Všechny mé objekty jsou nyní na dohodnutých cílových pozicích: P:(1,1), O:(3,1), T:(1,2), S:(2,2), A:(1,3), D:(7,6), M:(3,3). Jaký je tvůj stav? Jsou i tvé objekty na místech? Můžeme ukončit hru?"} \\
\\ 
\\ 
\texttt{Jaký je váš další příkaz?} \\
            }
        }
    }
    & \\ \\

    \theutterance \stepcounter{utterance}  
    & \multicolumn{4}{p{0.6\linewidth}}{
        \cellcolor[rgb]{0.9,0.9,0.9}{
            \makecell[{{p{\linewidth}}}]{
                \texttt{\tiny{[P1$\rangle$GM]}}
                \texttt{ŘEKNI: hotovo!} \\
            }
        }
    }
    & & \\ \\

    \theutterance \stepcounter{utterance}  
    & & & \multicolumn{2}{p{0.3\linewidth}}{
        \cellcolor[rgb]{0.9,0.9,0.9}{
            \makecell[{{p{\linewidth}}}]{
                \texttt{\tiny{[GM$|$GM]}}
                \texttt{ŘEKNI: hotovo!} \\
            }
        }
    }
    & & \\ \\

    \theutterance \stepcounter{utterance}  
    & & \multicolumn{4}{p{0.6\linewidth}}{
        \cellcolor[rgb]{0.9,0.9,0.9}{
            \makecell[{{p{\linewidth}}}]{
                \texttt{\tiny{[GM$\rangle$P2]}}
                \texttt{Vaše zpráva byla předána druhému hráči.} \\
\\ 
\\ 
\texttt{Právě hrajete kolo 10 z maximálně 28.} \\
\texttt{Dohromady jste nasbírali 2 z 16 pokut.} \\
\texttt{Druhý hráč poslal tuto zprávu:} \\
\\ 
\\ 
\texttt{"hotovo!"} \\
\\ 
\\ 
\texttt{Jaký je váš další příkaz?} \\
            }
        }
    }
    & \\ \\

    \theutterance \stepcounter{utterance}  
    & & & \multicolumn{4}{p{0.6\linewidth}}{
        \cellcolor[rgb]{0.9,0.9,0.9}{
            \makecell[{{p{\linewidth}}}]{
                \texttt{\tiny{[GM$\langle$P2]}}
                \texttt{ŘEKNI: hotovo!} \\
            }
        }
    }
     \\ \\

    \theutterance \stepcounter{utterance}  
    & & & \multicolumn{2}{p{0.3\linewidth}}{
        \cellcolor[rgb]{0.9,0.9,0.9}{
            \makecell[{{p{\linewidth}}}]{
                \texttt{\tiny{[GM$|$GM]}}
                \texttt{ŘEKNI: hotovo!} \\
            }
        }
    }
    & & \\ \\

    \theutterance \stepcounter{utterance}  
    & & & \multicolumn{2}{p{0.3\linewidth}}{
        \cellcolor[rgb]{0.9,0.9,0.9}{
            \makecell[{{p{\linewidth}}}]{
                \texttt{\tiny{[GM$|$GM]}}
                \texttt{true} \\
            }
        }
    }
    & & \\ \\

    \theutterance \stepcounter{utterance}  
    & & & \multicolumn{2}{p{0.3\linewidth}}{
        \cellcolor[rgb]{0.9,0.9,0.9}{
            \makecell[{{p{\linewidth}}}]{
                \texttt{\tiny{[GM$|$GM]}}
                \texttt{Initial states:} \\
\texttt{```} \\
\texttt{Player 1:} \\
\texttt{ 1234567} \\
\texttt{╔════╤══╗} \\
\texttt{║◌◌◌O│PM║ 1} \\
\texttt{║◌◌T◌├──╢ 2} \\
\texttt{║◌◌◌◌│◌◌║ 3} \\
\texttt{║S┌──┼──╢ 4} \\
\texttt{║◌│◌A│◌◌║ 5} \\
\texttt{╟─┴┬─┤◌◌║ 6} \\
\texttt{║◌◌│◌│D◌║ 7} \\
\texttt{╚══╧═╧══╝} \\
\\ 
\\ 
\texttt{Player 2:} \\
\texttt{ 1234567} \\
\texttt{╔════╤══╗} \\
\texttt{║◌◌◌◌│O◌║ 1} \\
\texttt{║P◌M◌├──╢ 2} \\
\texttt{║◌S◌◌│T◌║ 3} \\
\texttt{║◌┌──┼──╢ 4} \\
\texttt{║A│◌◌│◌◌║ 5} \\
\texttt{╟─┴┬─┤◌D║ 6} \\
\texttt{║◌◌│◌│◌◌║ 7} \\
\texttt{╚══╧═╧══╝} \\
\\ 
\texttt{```} \\
            }
        }
    }
    & & \\ \\

    \theutterance \stepcounter{utterance}  
    & & & \multicolumn{2}{p{0.3\linewidth}}{
        \cellcolor[rgb]{0.9,0.9,0.9}{
            \makecell[{{p{\linewidth}}}]{
                \texttt{\tiny{[GM$|$GM]}}
                \texttt{End states:} \\
\texttt{```} \\
\texttt{Player 1:} \\
\texttt{ 1234567} \\
\texttt{╔════╤══╗} \\
\texttt{║P◌O◌│◌◌║ 1} \\
\texttt{║TS◌◌├──╢ 2} \\
\texttt{║A◌M◌│◌◌║ 3} \\
\texttt{║◌┌──┼──╢ 4} \\
\texttt{║◌│◌◌│◌◌║ 5} \\
\texttt{╟─┴┬─┤◌D║ 6} \\
\texttt{║◌◌│◌│◌◌║ 7} \\
\texttt{╚══╧═╧══╝} \\
\\ 
\\ 
\texttt{Player 2:} \\
\texttt{ 1234567} \\
\texttt{╔════╤══╗} \\
\texttt{║P◌O◌│◌◌║ 1} \\
\texttt{║TS◌◌├──╢ 2} \\
\texttt{║A◌M◌│◌◌║ 3} \\
\texttt{║◌┌──┼──╢ 4} \\
\texttt{║◌│◌◌│◌◌║ 5} \\
\texttt{╟─┴┬─┤◌D║ 6} \\
\texttt{║◌◌│◌│◌◌║ 7} \\
\texttt{╚══╧═╧══╝} \\
\\ 
\texttt{```} \\
            }
        }
    }
    & & \\ \\

    \theutterance \stepcounter{utterance}  
    & & & \multicolumn{2}{p{0.3\linewidth}}{
        \cellcolor[rgb]{0.9,0.9,0.9}{
            \makecell[{{p{\linewidth}}}]{
                \texttt{\tiny{[GM$|$GM]}}
                \texttt{* success: True} \\
\texttt{* lose: False} \\
\texttt{* aborted: False} \\
\texttt{{-}{-}{-}{-}{-}{-}{-}} \\
\texttt{* Shifts: 6.00} \\
\texttt{* Max Shifts: 12.00} \\
\texttt{* Min Shifts: 6.00} \\
\texttt{* End Distance Sum: 0.00} \\
\texttt{* Init Distance Sum: 20.21} \\
\texttt{* Expected Distance Sum: 29.33} \\
\texttt{* Penalties: 2.00} \\
\texttt{* Max Penalties: 16.00} \\
\texttt{* Rounds: 10.00} \\
\texttt{* Max Rounds: 28.00} \\
\texttt{* Object Count: 7.00} \\
            }
        }
    }
    & & \\ \\

\end{supertabular}
}

\end{document}
